\DTMsavedate{mydate}{2022-06-30}
\maketitle{Peace, Leadership and Conflict Transformation \Romannum{1}}{Business Management}{\DTMusedate{mydate}}

% ----------------------------------------------------- %
% COURSE INFO
\subsection*{Course Information}\label{course:plc1101}
\vspace{0ex}%
\noindent\parbox{0.5\textwidth}{%
    \noindent\begin{tabular}{@{}ll}
         \textsf{Course:}          & 	Peace, Leadership and Conflict Transformation \Romannum{1}    \\
         \textsf{Course Code:}         & 	PLC 1101    \\
         \textsf{Credit Hours:}    & 	5
    \end{tabular}}
\vspace{2ex}\hrule\vspace{2ex}

% ----------------------------------------------------- %
% INSTRUCTOR INFO
\subsection*{Lecturer Information}
\noindent\parbox{0.5\textwidth}{%
    \noindent\begin{tabular}{@{}ll}
                 \textsf{Name:}         & David Foya \\
% \textsf{Phone:}        &    \phone          \\
\textsf{Email:}        &    \href{mailto:david.foya@nust.ac.zw}{david.foya@nust.ac.zw} \\
\textsf{Qualifications:}        & Phd in Peace and Governance, Master in Peace and Governance, BA in Education
    \end{tabular}}
\vspace{2ex}\hrule\vspace{2ex}


% ----------------------------------------------------- %
% COURSE DESCRIPTION
\subsection*{Course Synopsis}
The Peace, Leadership and Conflict Transformation course is tailored in a manner to provide
students with intellectual skills on the symbiotic relationship that exist on the three tier terms
[peace, leadership and conflict]. The course attempts to probe into the interplay between these
thematic motifs and show their role and complementarities in the process of human development.

The course further seeks to provide a skills tool kit on how to analyze conflicts, identify their
underlying causes, evaluate how conflict undermines the productive use of resources thereby
plaguing development and how responsible leadership transforms adversity into peaceful,
equitable and just global society in harmony with nature.

It is envisaged that the students who would successfully completed the course will be well
grounded in the theory and practice to face the challenges of leadership and conflict at personal,
community, national and global levels. The students would be able to trace the emerging patterns
and conflict trends in Africa shall form the basis of reflection.
\vspace{2ex} \hrule\vspace{2ex}

% ----------------------------------------------------- %
% LEARNING OBJECTIVES
\subsection*{Learning Objectives}
The objectives of this course are to provoke debate , through critical thinking, about conflict and
conflict management, focusing on both endogenous and exogenous styles of leadership and
conflict resolution approaches(negotiation, mediation, reconciliation etc). Therefore students
should be able to:
\begin{outline}
    \1      Describe and contextualise key theories and concepts of peace, leadership and conflict.

    \1  	Examine the origins and dynamism of both leadership and conflict.

    \1  	Understand the art of and functions of leadership in society and organisations and appreciate the interplay between leadership and conflict management.

    \1  	Develop descriptive and analytical research skills in leadership and conflict transformation.

\end{outline}
\vspace{2ex}\hrule\vspace{2ex}



% ----------------------------------------------------- %
% LEARNING OUTCOMES
\subsection*{Learning Outcomes}
By the end of this course, participants will be able to:

\begin{outline}
    \1      Demonstrate use of acquired skills by applying them to their context and able to transform the lives of the members of the community.

    \1  	Have a verbal or written demonstration of learnt concepts.

    \1  	Be able to explain the rationale for carrying out research in the three areas of peace, leadership and conflict transformation.

    \1  	Be able to identify research area, conduct research, gather data, analyse, interpret and make a write up.

\end{outline}
\vspace{2ex}\hrule\vspace{2ex}

% ----------------------------------------------------- %
% LEARNING OUTCOMES
\subsection*{Learning Requirements}
\begin{outline}
    \1      Students are expected to attend classes regularly. Attendance of classes is a requirement and excessive absence will result in failing the course.

    \1  	Written assignments should be typed, double spaced, (12 font size), with reference/bibliography at the end.

    \1  	Students are expected to participate in class discussions. Reading course materials in advance is essential for participation in discussion. Oral presentations shall include reviews, case studies from field visits, descriptive and analytical reading of the selected course materials.

    \1  	Written assignments should meet the university general requirements.

    \1  	Any written assignments should provide empirical description as evidence to substantiate on claims.

\end{outline}
\vspace{2ex}\hrule\vspace{2ex}



% ----------------------------------------------------- %
% COURSE TOPICS
\subsection*{Topics}
\begin{outline}
    \1 Concept of Conflict and Peace
    \1 Theories of Conflict and Peace
    \1 Theories of Conflict
    \1 Classical Social Structural Theories of Conflict
    \1 Modern Structural Theories of Conflict
    \1 Cultural, religious, Ethnic and Identity Based Conflicts
    \1 Gender and Conflict
    \1 Impact of Conflict on development
    \1 Conflict Resolution Process

\end{outline}
\vspace{2ex} \hrule\vspace{2ex}

% ----------------------------------------------------- %
% Students Assenssment
\subsection*{Assessment of Students}
\begin{outline}
    \1 Students are assessed through written assignments and tests that shall form part of course work assessment mark.
    \1 Written assignments should meet University general requirements.
    \1 Students are expected to participate in class discussion and to read course material.
    \1 There shall be a three-hour final examination at the end of each semester.
    \1 Marking of scripts rewarding of marks shall be based on:
        \begin{enumerate}
            \item logic, analytical rigour and coherence.
            \item understanding of the issues at hand.
            \item creativity and critical engagement with the topic.
            \item style, suave and presentation.
            \item in the case of assignment, bibliography should be given.
        \end{enumerate}

\end{outline}
\vspace{2ex}\hrule\vspace{2ex}

% ----------------------------------------------------- %
% Grade Evaluation
\subsection*{Course Evaluation and Grading Policies}
Grades are based on:
    Continuous assessment (\qty{30}{\percent}),
    Final Exam (\qty{70}{\percent})
\vspace{2ex}\hrule\vspace{2ex}

% ----------------------------------------------------- %
% EQUIPMENT
% https://sites.udel.edu/computing-purchases/personal-specs/
% \subsection*{Essential Equipment}

% \vspace{2ex}\hrule\vspace{2ex}

% Policies and Other information
\subsection*{Policies and Other Information}
% Conduct
\subsection*{Key Text}
\begin{itemize}
    \item Aye, J (2000) Leadership in Contemporary Africa: Exploratory Study, New York United Nations.
    \item Banks, M. Four Conceptions of Peace: Unpublished Article.
    \item Deng , Z. (2002) Developmental Challenges in Africa. London, OUP.
    \item Galtung, J. (1997). Peace by Peaceful Means. London: Sage Press.
    \item Lederach, J.P. (1995). Preparing for Peace: Conflict Transformation Across Cultures. Syracuse: Syracuse University Press.
    \item \href{www.beyondintractability.org}{www.beyondintractability.org}
    \item \href{www.berghof-handbook.net}{www.berghof-handbook.net}
\end{itemize}
\vspace{2ex}\hrule
\vspace{2ex}\hrule\vspace{2ex}